% Options for packages loaded elsewhere
\PassOptionsToPackage{unicode}{hyperref}
\PassOptionsToPackage{hyphens}{url}
\documentclass[
]{article}
\usepackage{xcolor}
\usepackage{amsmath,amssymb}
\setcounter{secnumdepth}{-\maxdimen} % remove section numbering
\usepackage{iftex}
\ifPDFTeX
  \usepackage[T1]{fontenc}
  \usepackage[utf8]{inputenc}
  \usepackage{textcomp} % provide euro and other symbols
\else % if luatex or xetex
  \usepackage{unicode-math} % this also loads fontspec
  \defaultfontfeatures{Scale=MatchLowercase}
  \defaultfontfeatures[\rmfamily]{Ligatures=TeX,Scale=1}
\fi
\usepackage{lmodern}
\ifPDFTeX\else
  % xetex/luatex font selection
\fi
% Use upquote if available, for straight quotes in verbatim environments
\IfFileExists{upquote.sty}{\usepackage{upquote}}{}
\IfFileExists{microtype.sty}{% use microtype if available
  \usepackage[]{microtype}
  \UseMicrotypeSet[protrusion]{basicmath} % disable protrusion for tt fonts
}{}
\makeatletter
\@ifundefined{KOMAClassName}{% if non-KOMA class
  \IfFileExists{parskip.sty}{%
    \usepackage{parskip}
  }{% else
    \setlength{\parindent}{0pt}
    \setlength{\parskip}{6pt plus 2pt minus 1pt}}
}{% if KOMA class
  \KOMAoptions{parskip=half}}
\makeatother
\usepackage{longtable,booktabs,array}
\usepackage{caption}
\captionsetup[table]{skip=6pt}
\newcounter{none} % for unnumbered tables
\usepackage{calc} % for calculating minipage widths
% Correct order of tables after \paragraph or \subparagraph
\usepackage{etoolbox}
\makeatletter
\patchcmd\longtable{\par}{\if@noskipsec\mbox{}\fi\par}{}{}
\makeatother
% Allow footnotes in longtable head/foot
\IfFileExists{footnotehyper.sty}{\usepackage{footnotehyper}}{\usepackage{footnote}}
\makesavenoteenv{longtable}
\setlength{\emergencystretch}{3em} % prevent overfull lines
\providecommand{\tightlist}{%
  \setlength{\itemsep}{0pt}\setlength{\parskip}{0pt}}
\usepackage{bookmark}
\IfFileExists{xurl.sty}{\usepackage{xurl}}{} % add URL line breaks if available
\urlstyle{same}
% fallback for those not using the hyperref driver hyperxmp:
\makeatletter
\@ifundefined{xmpquote}{\newcommand{\xmpquote}[1]{#1}}{}
\makeatother
\hypersetup{
  hidelinks,
  pdfcreator={LaTeX via pandoc}}

\author{}
\date{}

\begin{document}

\section{Gate E7 report --- run\_v1.2\_r1 (Paper IV
v1.2)}\label{gate-e7-report--run_v12_r1-paper-iv-v12}

Verdict: \textbf{PASS with MINOR findings}. Evidence directory:
\texttt{20\_EVIDENCE/E7/}.

\begin{center}\rule{0.5\linewidth}{0.5pt}\end{center}

\subsection{E7 (attribution / literature) --- first
pass}\label{e7-attribution--literature--first-pass}

Target:
\texttt{01\_manuscript/v1.2\_full\_rebuild\_candidate/PAPER\_IV\_preprint\_v1.2\_en.md},
sha256
\texttt{a501d53b527da22f20112db2b05917fb61f46dd8e080faf14b69349b240cc152}
(2464 lines). Audit date 2026-09-30. Labels: \textbf{SV} =
SOURCE\_VERIFIED (read in the saved source), \textbf{AI} =
AUDITOR\_INFERENCE, \textbf{UV} = UNVERIFIED. Line numbers refer to the
manuscript unless prefixed by a source. Page numbers for {[}15{]} are
those printed in arXiv:2609.20871v1. Quotes from sources are kept under
15 words; everything else is paraphrase.

\subsection{0. Sources and later-version
status}\label{0-sources-and-later-version-status}

{\def\LTcaptype{none} % do not increment counter
\begin{longtable}[]{@{}lllll@{}}
\toprule\noalign{}
Ref & URL & Version retrieved & Source date & sha256 (saved file) \\
\midrule\noalign{}
\endhead
\bottomrule\noalign{}
\endlastfoot
{[}15{]} & \url{https://arxiv.org/pdf/2609.20871v1} & v1 (cited) & 15
Sep 2026 & \texttt{9654af66…d5fa53} (okechukwu\_v1.pdf) \\
{[}15{]} abs & \url{https://arxiv.org/abs/2609.20871} & abs page & --- &
\texttt{e1ce80ea…b1b0} \\
{[}22{]} & \url{https://arxiv.org/pdf/1902.06135v1} & v1 (cited) & 16
Feb 2019 & \texttt{df77c1d9…6934} (verclos\_v1.pdf) \\
{[}22{]} abs & \url{https://arxiv.org/abs/1902.06135} & abs page & --- &
\texttt{af8ac080…a1b0} \\
{[}5{]} PDF & raw.githubusercontent
\ldots/cbde8a0a\ldots/manuscript/main.pdf & commit cbde8a0a & printed 8
Sep 2026 & \texttt{53e15161…3930} \\
{[}5{]} README / FORMALIZATION\_STATUS & raw
\ldots/cbde8a0a\ldots/README.md, lean/FORMALIZATION\_STATUS.md & commit
cbde8a0a & 2026-09-08 & \texttt{0e08ab41…0b30b} /
\texttt{05cceb10…93f4} \\
{[}8{]} & \url{https://www.erdosproblems.com/81} & live, 2026-09-30 &
page last edited 28 Dec 2025 & \texttt{a02ca0e2…2543} \\
{[}21{]} &
\url{https://www.erdosproblems.com/forum/thread/81/proof-claims} & live,
2026-09-30 & claim 201 dated 2026-08-10 & \texttt{69d52613…b329} \\
{[}23{]} & Crossref for 10.1137/06064888X & metadata & 2008 &
\texttt{08dfd5ad…87cb} \\
\end{longtable}
}

Full list with URLs and retrieval times: \texttt{SOURCES\_SHA256.txt}.
Queries: \texttt{SEARCH\_LOG.md}.

Later-version status as of 2026-09-30 (SV):

\begin{itemize}
\tightlist
\item
  {[}15{]}: the arXiv submission history lists only v1 (15 Sep 2026
  16:25:32 UTC), 24 pages, with no journal reference.
\item
  {[}22{]}: the history lists only v1 (16 Feb 2019), with no journal
  reference or DOI. Crossref and web search found no journal version (UV
  that none exists; the dblp query failed).
\item
  {[}5{]}: \texttt{main} HEAD is still \texttt{cbde8a0a} (committed
  2026-09-08T10:47:29Z; repo pushed\_at 2026-09-08T12:53:39Z). Tag
  \texttt{v0.1.0-proof-claim} points to the same commit. The only other
  branch is \texttt{codex/revise-manuscript-lean-audit} at 9944194c, an
  ancestor merged via PR \#1. \texttt{manuscript/main.pdf} has the same
  sha256 at the cited commit, at \texttt{main} and in the release asset
  digest. \textbf{Nothing has changed since the cited commit.}
\item
  {[}8{]}: Problem \#81 is still marked \textbf{OPEN} on the site.
\end{itemize}

\subsection{A. Theorem C (l.82--97) vs {[}15, Theorem 1.1{]} ---
component
table}\label{a-theorem-c-l8297-vs-15-theorem-11--component-table}

{\def\LTcaptype{none} % do not increment counter
\begin{longtable}[]{@{}llll@{}}
\toprule\noalign{}
Component & {[}15{]} v1 (SV, p.1--2, 5--6, 19--20) & Manuscript v1.2 &
Verdict \\
\midrule\noalign{}
\endhead
\bottomrule\noalign{}
\endlastfoot
Graph class & s-simplicial: d(v) − ω(N(v)) ≤ s. rsd(G) is the least s
such that, in every induced H and outside every proper clique of H,
there is an s-simplicial vertex. "The empty clique is allowed." (p.2) &
l.84: in every induced U, outside every prescribed clique R ⊊ U, some
vertex\textquotesingle s neighbourhood contains a clique omitting ≤ s
neighbours. The empty root is allowed and the definition is credited to
{[}15, §1{]}. & OK (equivalent formulation, attributed) \\
Target & F\_n = ⌊n(n+1)/6⌋, Q\_s(n) = F\_\{n+s\} − C(s+1,2) (1.1) &
Q\_s(n) = M(n+s) − C(s+1,2), with notation credited to {[}15,(1.1){]}
(l.84, l.90) & OK \\
Quantifiers & ∀s ∃K\_s ≥ 0, N\_s. The bound cp ≤ Q\_s + K\_s holds at
\textbf{all} orders; for n ≥ N\_s the maximum is exactly Q\_s(n). & ∀s
∃N\_s ∀n ≥ N\_s: cp ≤ c\_4 ≤ Q\_s(n). Eventual only. & OK. The eventual
scope is stated; the all-orders clause is covered in §8.1 (see G). \\
Piece size & Statement is for cp only. The explicit constructions (Lemma
2.5, Lemmas 4.1--4.3, (5.11)) use only edges and triangles. Lemma 4.2 is
stated for q\_3. & c\_4 (pieces of order ≤ 4) & See A.2. The stated
strengthening is literally correct. \\
Attainment & Lemma 2.5, at every sufficiently large order & l.93:
attainment for each sufficiently large order, also for unrestricted cp &
OK. Credited at l.97. \\
Extremal classification & Equality holds exactly for (K\_\{k−s\} ⊔ K̄\_s)
∨ K̄\_\{n−k\}, with k nearest to (2(n+s)+1)/6 (1.3). Complement bars
confirmed visually from the rendered p.2. & (6.18): cp = Q\_s ⟺ c\_4 =
Q\_s ⟺ G ≅ (K\_\{k−s\} ⊔ I\_s) ∨ I\_\{n−k\}, same k (l.1001--1004) & OK.
Same family, credited at l.97 and l.1014. \\
Witness family & book graphs (K\_\{k−s\} ⊔ K̄\_s) ∨ K̄\_b & same (Appendix
E witness) & OK \\
Constants & K\_s, N\_s existential & N\_s explicit (tower, Cor. 6.6) &
OK, no overclaim (l.1126) \\
\end{longtable}
}

\subsubsection{A.1 Role of the cutoff L in {[}15{]}
(SV)}\label{a1-role-of-the-cutoff-l-in-15-sv}

\begin{itemize}
\tightlist
\item
  \textbf{§3 (p.6).} L ≥ 3 is fixed and q\_L and q\_L\^{}* are defined.
  Lemma 3.1 requires L ≥ 4. Its final assertion is q\_L\^{}* ≤ M\_n +
  n·rsd(G).
\item
  \textbf{Theorem 3.3 (p.8--9).} For every ε, it produces an integer L ≥
  4 and constants. In the proof, ξ is chosen first and then L, so that
  33(ξ + 1/{[}2(L−1){]}) \textless{} ε/4. Hence L grows as ε → 0.
\item
  \textbf{Lemma 3.4 (p.10).} For every fixed L and ζ \textgreater{} 0,
  q\_L ≤ q\_L\^{}* + ζ\textbar G\textbar² for all sufficiently large
  graphs. The proof bundles templates J and uses Haxell--Rödl/Yuster.
\item
  \textbf{Proof of Theorem 1.3 (p.11).} L, δ\_0, η\_0 are taken from
  Theorem 3.3 ("The cutoff L is now fixed."). Lemma 3.4 is then applied
  with error η\_0 n²/2. Its role is to turn a lower bound on cp (through
  q\_L ≥ cp) into a lower bound on q\_L\^{}*.
\item
  \textbf{Proposition 5.4 (p.17) and Theorem 1.1 (p.19).} Prop. 5.4
  calls Theorem 1.3 to obtain the cliques needed by Lemma 5.1. Theorem
  1.1 uses Prop. 5.4 in a smallest-counterexample argument. L appears
  nowhere else.
\item
  \textbf{(1.5), p.11.} The source says to "use L = 4 in Lemmas 3.1 and
  3.4". This is the only place where L = 4 is chosen.
\item
  \textbf{Verdict on l.1255.} The statement is \textbf{accurate} (SV):
  Theorem 3.3 chooses L from the accuracy; Prop. 5.4 and Theorem 1.1
  reach it through Theorem 1.3 and Lemma 3.4; and L = 4 is used only for
  (1.5).
\end{itemize}

\subsubsection{A.2 Is the order-four strengthening trivially available
from {[}15{]}?
(AI)}\label{a2-is-the-order-four-strengthening-trivially-available-from-15-ai}

\begin{itemize}
\tightlist
\item
  \textbf{What {[}15{]} proves literally.} It proves cp ≤ Q\_s(n) for n
  ≥ N\_s and states no bound on c\_4 or q\_L for non-extremal graphs. In
  its near-extremal regime (hypotheses of Prop. 5.4: δ(G) ≥ n/3 − o(n)
  and cp ≥ Q\_s), the construction (5.11) goes through Lemma 4.2, whose
  conclusion is a q\_3 bound. Graphs with cp \textless{} Q\_s(n) get no
  constructive treatment: the smallest-counterexample argument only
  needs the value of cp.
\item
  \textbf{Why simple substitution fails.} Running the §5.4 argument with
  c\_4 in place of cp would need Theorem 1.3 under a c\_4 (= q\_4)
  lower-bound hypothesis. Lemma 3.4 with L = 4 then bounds q\_4\^{}*
  from below. Theorem 3.3, however, needs a lower bound on q\_L\^{}* for
  a large L = L(ε), and q\_L\^{}* ≤ q\_4\^{}* goes the wrong way. So the
  literal chain does not give c\_4 ≤ Q\_s by substitution.
\item
  \textbf{Limit of this finding.} This is a statement about the literal
  chain only, not a claim that {[}15{]}\textquotesingle s method cannot
  be adapted. The manuscript makes the same disclaimer at l.1245 and
  l.1255.
\item
  \textbf{Verdict.} Calling Theorem C an order-four strengthening of the
  \emph{stated} bound (l.95, l.1255) is \textbf{OK}.
\end{itemize}

\subsubsection{\texorpdfstring{A.3 Mechanism description of {[}15{]} in
Table 4 / l.1210 and l.1257 --- \textbf{MINOR} (AI, grounded in SV
reading)}{A.3 Mechanism description of {[}15{]} in Table 4 / l.1210 and l.1257 --- MINOR (AI, grounded in SV reading)}}\label{a3-mechanism-description-of-15-in-table-4--l1210-and-l1257--minor-ai-grounded-in-sv-reading}

\begin{itemize}
\tightlist
\item
  \textbf{What the manuscript says.} Table 4 gives
  {[}15{]}\textquotesingle s positive-gain pieces as "K\_3,\ldots,K\_L
  in the far transfer", with "cp, via q\_L in that branch". Line 1257
  says Lemma 3.4 "pays the budget when ζn² ≤ M(n) − q\_L\^{}*(G)" and
  "is used in the branch with the corresponding margin".
\item
  \textbf{What {[}15{]} actually does.} The proof of Theorem 1.1 has no
  far-regime branch in which a margin pays for a transfer loss. Lemma
  3.4 enters only inside localization (proof of Thm 1.3), where it
  converts cp ≥ n²/6 − o(n²) into a lower bound on q\_L\^{}*. The one
  place it serves as an upper bound is (1.5), with L = 4. The explicit
  partitions {[}15{]} builds for Theorem 1.1 use only K\_2 and K\_3.
\item
  \textbf{Why it matters.} The table therefore gives {[}15{]} an
  architecture it does not have. A reader could infer that
  {[}15{]}\textquotesingle s partitions use large cliques. This does not
  affect priority claims, and it arguably understates {[}15{]}.
\item
  \textbf{Recommendation.} Rewrite Table 4\textquotesingle s {[}15{]}
  row and l.1257 to say that Lemma 3.4 is used for localization and for
  (1.5). Also say that the near-extremal constructions in {[}15{]} are
  edge/triangle constructions.
\end{itemize}

\subsection{B. Checks of cited items of
{[}15{]}}\label{b-checks-of-cited-items-of-15}

{\def\LTcaptype{none} % do not increment counter
\begin{longtable}[]{@{}lll@{}}
\toprule\noalign{}
Manuscript claim & {[}15{]} v1 (SV) & Verdict \\
\midrule\noalign{}
\endhead
\bottomrule\noalign{}
\endlastfoot
{[}15, Cor 1.2{]} is the chordal case (l.43); optimal cores nearest to
(2n+1)/6 (l.1263) & Cor 1.2, p.2: cp ≤ F\_n + O(1); eventual max F\_n;
equality exactly K\_k ∨ K̄\_\{n−k\}, k nearest to (2n+1)/6 & OK \\
{[}15, Thm 1.3{]} structural sequence result (l.43, l.1261) & Thm 1.3,
p.4: rsd = o(n), cp ≥ n²/6 − o(n²) give cliques A\_j with ‖A\_j‖ = n/3 +
o(n) and D\_A + e(G−A) = o(n²); plus (1.5) & OK \\
{[}15, Cor 5.5{]} qualitative edit stability (l.43, l.1259) & Cor 5.5,
p.20: for fixed s and ε, there are δ, N such that cp ≥ Q\_s − δn² gives
≤ εn² edits to form (1.3). Qualitative. & OK \\
{[}15, Lemma 3.1{]} signed-star argument (l.45, l.1070, F.2 l.2150) &
Lemma 3.1, p.7 (§3.1 "A maximal positive clique-star"). The
manuscript\textquotesingle s (F.2) rows, the α = 5c/12 split and the
envelope c(5n−4c−1)/12 match {[}15{]} p.7--8. The manuscript adds the
capped exception charge Σ min(s,j). & OK, adequately attributed ("cited
as a method", l.1070) \\
{[}15, Lemma 3.4{]} q\_L ≤ q\_L\^{}* + ζn² (l.1257); exact fractional
partition ⟺ gain model & Lemma 3.4, p.10, including the "equivalently a
fractional packing" remark in its proof & OK for the statement;
mechanism framing MINOR (A.3) \\
{[}15, Theorem 3.3{]} n/6 + 1/24 term (l.1263) & p.9: ϱ ≤ (δ\_0+η\_0)n²
+ n/6 + 1/24 & OK \\
{[}15, §7{]} quantitative version of Cor 5.5 (l.1259) & §7, p.23: asks
for the optimal dependence of edits on Q\_s(n) − cp(G), n and s. Notes
that the proof gives only qualitative stability. & OK \\
Table 6 (l.1249--1251): signed dual Thm 3.3; Haxell--Rödl + joint
templates Lemma 3.4; list colouring + exceptions §§4--5 & Thm 3.3 uses
the signed dual; Lemma 3.4 uses bundled templates J with
Haxell--Rödl/Yuster; §4 uses Galvin\textquotesingle s
list-edge-colouring theorem and §5 a bounded exceptional set & OK \\
Order-three obstruction (Prop D.2, l.95, l.1748--1761) & Neither
{[}15{]} nor {[}5{]} makes any claim about global order-3 sufficiency or
insufficiency. {[}15{]}\textquotesingle s near-extremal constructions
are order ≤ 3 (Lemma 4.2 is a q\_3 bound). {[}5{]} builds a K\_2/K\_3
partition near the extremal family ((5.8): cp≤3 ≤ M(n)) and needs K\_4
only through the fractional far functional (p.2 remark on complete
graphs). & OK. No contradiction. The manuscript restricts the claim to s
= 0 and uniformity (l.95). \\
\end{longtable}
}

\subsection{C. Theorem C′ / Cor 6.3 / Prop 6.3a / Prop 6.4 / Thm 6.5 vs
{[}15{]} --- component
table}\label{c-theorem-c--cor-63--prop-63a--prop-64--thm-65-vs-15--component-table}

{\def\LTcaptype{none} % do not increment counter
\begin{longtable}[]{@{}llll@{}}
\toprule\noalign{}
Aspect & {[}15{]} & Manuscript & Verdict \\
\midrule\noalign{}
\endhead
\bottomrule\noalign{}
\endlastfoot
Defect regime & Cor 5.5: fixed s. Thm 1.3: rsd = o(n). & C′ / Cor 6.3:
fixed s. Thm 6.5: s\_j = o(n\_j). & Parallel structure, correctly
matched \\
Near-extremality hypothesis & on cp & on c\_4 (c\_4 ≥ Q\_s − δ). This is
implied by the cp hypothesis, so it is a \textbf{weaker} hypothesis and
a stronger theorem. & See C.1 \\
Quantitative vs qualitative & Cor 5.5 qualitative (ε, δ); §7 explicitly
says there is no rate & d\_E ≤ A\_sδ (non-optimal core) and d\_E ≤ A\_sδ
+ n√((1+4A\_s)δ) (optimal size, 6.17); explicit γ\_s, N\_s\^{}stab &
Positioning as a "quantitative form" of Cor 5.5 (l.43, which explicitly
includes Cor 6.3) is accurate \\
Square-root lower scale & none & Prop 6.3a family & No conflict with
{[}15{]} \\
Partition control & {[}15{]} contains \textbf{no} statement about the
structure of near-optimal partitions (whole text read) & (6.13), (6.26):
one root controls all partitions & "adds simultaneous partition control"
(l.43, l.1261) accurate \\
Linear / quantitative stability in {[}15{]}? & None. The only
quantitative objects are Lemma 3.1 / (1.5) and Thm 6.2 (finite signed
bound). & --- & Manuscript does not misstate {[}15{]} here \\
Thm 1.3 vs Thm 6.5 & cliques A\_j, ‖A\_j‖ = n/3 + o(n), D\_A + e(G−A) =
o(n²) & R\_j, ‖R\_j‖ = n/3 + o(n), d\_E(G\_j, S\_\{R\_j\}) = o(n²), with
weaker c\_4 hypothesis & "recovers the structural conclusion"
accurate \\
(1.5) vs Prop 6.4 + (6.24) & Stated for cp: (2n+1)²/24 + n·rsd + ϵ(n)n²
for every n ≥ 8, with ϵ(n) → 0 graph-independent. Its proof uses L = 4
in Lemmas 3.1 and 3.4, so it proves a q\_4 = c\_4 bound. & c\_4 ≤ M(n−s)
+ ns (−C(s+1,2) on 4s ≤ n) + εn², with N(ε) independent of G and s.
Leading s-term ≈ 2ns/3 vs ns. & See C.2 \\
\end{longtable}
}

\subsubsection{\texorpdfstring{C.1 Wording error at l.1259 ---
\textbf{MINOR}
(AI)}{C.1 Wording error at l.1259 --- MINOR (AI)}}\label{c1-wording-error-at-l1259--minor-ai}

\begin{itemize}
\tightlist
\item
  \textbf{The wording.} Line 1259 calls the c\_4 near-extremality
  hypothesis "the stronger near-extremality hypothesis on c\_4". It then
  says the result also applies when the hypothesis is stated for cp,
  since cp ≤ c\_4.
\item
  \textbf{Why it is wrong.} The hypothesis c\_4 ≥ Q\_s − δ is implied by
  cp ≥ Q\_s − δ, so it is the weaker hypothesis. The sentence
  contradicts itself.
\item
  \textbf{Fix.} Replace "stronger" by "weaker" or "c\_4-form of the".
\end{itemize}

\subsubsection{\texorpdfstring{C.2 "order-four version of its
approximation (1.5)" (l.1261) --- \textbf{MINOR} (AI grounded in SV
p.11)}{C.2 "order-four version of its approximation (1.5)" (l.1261) --- MINOR (AI grounded in SV p.11)}}\label{c2-order-four-version-of-its-approximation-15-l1261--minor-ai-grounded-in-sv-p11}

\begin{itemize}
\tightlist
\item
  \textbf{Why the phrase overstates.} {[}15{]} states (1.5) for cp, but
  its one-line proof already uses L = 4, so it literally yields a q\_4
  bound. The manuscript itself acknowledges this at l.1255.
\item
  \textbf{Fix.} Describe the contribution as a refined finite-order and
  s-dependence form of (1.5), noting that {[}15{]}\textquotesingle s
  proof of (1.5) already works with L = 4.
\end{itemize}

\subsection{D. {[}22{]} de Joannis de Verclos (SV unless
noted)}\label{d-22-de-joannis-de-verclos-sv-unless-noted}

\begin{itemize}
\item
  \textbf{What {[}22{]} proves.} Theorem 1 (p.2): chordality is testable
  with query complexity O(ε\^{}\{-37\}), i.e. polynomial ("easily
  testable"). The contrapositive is a polynomial removal-type statement.
  The manuscript\textquotesingle s wording "polynomial chordal-removal
  development" (l.1138) is fair (AI).
\item
  \textbf{Lemma 3 (p.3).} It assumes ε \textgreater{} 0 and n ≥
  ε\^{}\{-1\}, and concludes that G is 6ε\^{}\{1/2\}-close to chordal.
  The manuscript (F.10, l.2287) says constant 8 instead of 6, with no n
  ≥ 1/ε condition, and "neither uniformly stronger". This matches.
  \textbf{OK}
\item
  \textbf{Lemma 4 (p.3).} It is the Gyárfás--Hubenko--Solymosi
  large-clique lemma, cited within {[}22{]}. Line 2280 says the
  manuscript "replaces the use of {[}22, Lemma 4{]}" and "is not
  presented as a new version of the full theorem cited there".
  \textbf{OK.} The manuscript could name GHS explicitly; optional.
\item
  \textbf{Lemma 11 (p.10--14).} It concerns pinned chordal graphs, with
  T having at most k leaves and m11 = 2\^{}40 ε\^{}\{-4\} k\^{}6 ln\^{}4
  k. Its proof has three parts that match the
  manuscript\textquotesingle s three cases:

  \begin{itemize}
  \tightlist
  \item
    Claim 3: every colouring has ≥ δn² conflicts, so Theorem 2 (set
    colouring) applies. This is l11\_caseA.
  \item
    Claim 5: a section G{[}L\_i,R\_i{]} far from M2-free gives an
    induced C4 w.p. ≥ 1/2. This is l11\_caseB1 / claim5\_count.
  \item
    Claim 6 and the gluing construction: F chordal and
    \textbar E(F)△E(G)\textbar{} ≤ δn² + (ε/2)n² ≤ εn². This is
    l11\_caseB2 / (F.12).
  \end{itemize}

  The manuscript\textquotesingle s attribution ("retain their roles from
  the removal strategy of {[}22{]}"; no claim that the published Lemma
  11 has the modified statement, l.2341) is \textbf{accurate}.
\item
  \textbf{Table 10 (l.2291--2300).} Four claims were checked against
  {[}22{]}:

  \begin{itemize}
  \tightlist
  \item
    {[}22{]}\textquotesingle s pinned representation is on the fixed
    tree T (Definition 2, p.9), so "enlarged class" with a
    path-preserving map is a correct description of a change.
  \item
    Lemma 11\textquotesingle s Claim 1 (gate set of size ≤ 3k built from
    leaves) is correctly identified as "not imported".
  \item
    Theorem 2 has m2 = 36k ln(max\textbar L\_u\textbar)ε\^{}\{-2\},
    matching "published logarithmic bounds".
  \item
    The manuscript\textquotesingle s m11 = ⌈2\^{}58(K+1)\^{}15/ε\^{}12⌉
    is larger than {[}22{]}\textquotesingle s. The manuscript says so
    (l.2300).
  \end{itemize}

  \textbf{OK}
\item
  \textbf{Journal version.} None found (UV). The citation to arXiv v1 is
  correct and current.
\item
  \textbf{Overall verdict D: OK.} The description of {[}22{]} is
  accurate.
\end{itemize}

\subsection{E. {[}5{]} (N0zoM1z0/erdos-81 @
cbde8a0a)}\label{e-5-n0zom1z0erdos-81--cbde8a0a}

\begin{itemize}
\item
  \textbf{Eventual chordal maximum.} SV, {[}5{]} p.2 and p.15: Theorem
  1.1 (eventual max M(n)), Theorem 9.1 (cp≤4 ≤ M(n) for n ≥ N) and
  Corollary 1.2 (all orders). The manuscript\textquotesingle s l.43,
  l.61 and l.1265 are \textbf{OK}.
\item
  \textbf{FORMALIZATION\_STATUS.md.} It is dated 2026-09-08. Its
  theorems take \texttt{Erdos81.ExternalInputs.Inputs} with three
  fields: VizingInput, HaggkvistJanssenInput and PackingTransferInput.
  The last returns certified rational primal/dual optima and an attained
  integral optimum. These are stated to be explicit parameters, not
  axioms. The manuscript\textquotesingle s l.1132 is \textbf{accurate},
  including "does not assert that the mathematical proof of {[}5{]} is
  conditional on conjectures". Additional context: {[}5{]} attributes
  its transfer to Rohatgi--Urschel--Wellens {[}RUW21, Thm 3.4/3.6{]}
  (p.3).
\item
  \textbf{Table 5 section locations} (l.1230--1238, l.1217):

  \begin{itemize}
  \tightlist
  \item
    §2: external inputs / transfer statement
  \item
    §3: integral construction around a root (Vizing + Häggkvist--Janssen
    list colouring)
  \item
    §4: strict root regularization
  \item
    §5: quantitative local stability
  \item
    §6: monotone copy path
  \item
    §8: global stability by first entry
  \item
    §9: main theorem, cp≤4 ≤ M(n)
  \end{itemize}

  All \textbf{match} (SV headings).
\item
  \textbf{Strict root regularization (l.1267).} {[}5{]} Lemma 4.2 (p.7)
  gives cp≤3(G) ≤ p′q′ − C(p′,2) − m′/9 − A′/2 for a new clique P′. This
  equals B\_n(p′) − m′/9 − A′/2 with q′ = n − p′. \textbf{OK} (primes
  omitted in the manuscript; trivial).
\item
  \textbf{Credits to the series (l.1241).} {[}5{]}\textquotesingle s
  introduction (p.1--2) cites Paper II {[}TG26a{]} (fractional
  functional), Paper III {[}TG26b{]} (split linear-error bound) and
  Cipollini. §6 (p.9) says the copy inequality, clone-class lift and
  complete-split terminal reduction follow the strategy of {[}TG26a,
  §§3--5{]}. \textbf{OK.}
\item
  \textbf{Changes since cbde8a0a.} None (see §0).
\item
  \textbf{Attribution of {[}5{]} as "Anonymous" (l.2428) --- MINOR /
  disclosure note} (SV facts, AI assessment).

  \begin{itemize}
  \tightlist
  \item
    The printed PDF author is "ANONYMOUS", so the
    manuscript\textquotesingle s choice follows the document.
  \item
    Public attribution exists elsewhere, though. The README credits
    Morluto with the mathematical proof (GPT-assisted) and N0zoM1z0 with
    the manuscript and Lean. erdosproblems.com proof claim 285
    (2026-09-08) is by "Morluto, Yang Luo, Yinzhen Huang, Grace Lee" and
    links this repository.
  \item
    The manuscript\textquotesingle s own tool reference {[}14{]}
    (Jacobian, github.com/morluto/jacobian, l.2414, l.2446) belongs to
    the same Morluto. {[}5{]}\textquotesingle s README and claim 285 say
    Jacobian played a substantial role in deriving {[}5{]}.
  \item
    \textbf{Suggestions.} (i) Add a pointer to proof claim 285 in
    {[}5{]}\textquotesingle s reference. (ii) Consider a one-line
    disclosure that the tool {[}14{]} is by an author of the {[}5{]}
    claim. This is not an accusation of dependence; the
    manuscript\textquotesingle s ConeAudit (l.1136) concerns formal
    namespace exclusion only.
  \end{itemize}
\end{itemize}

\subsection{F. Cipollini {[}21{]}}\label{f-cipollini-21}

\begin{itemize}
\tightlist
\item
  \textbf{Anchor and date.} SV: anchor \texttt{\#proof-claim-201} exists
  and is labelled a "partial proof claimed by Ricky Cipollini". It was
  submitted 2026-08-10 17:06:33.
\item
  \textbf{Content.} The summary claims partition into n²/6 + o(n²)
  cliques and names the o(n²)-to-O(n) gap as the remaining difficulty.
  The external link is the same Overleaf ID \texttt{thjptfhgnmxc} (the
  site appends \texttt{\#cc1388}).
\item
  \textbf{Manuscript.} Lines 47 and 2460 are \textbf{accurate}. The
  Overleaf manuscript was not opened (coverage limit).
\end{itemize}

\subsection{G. Historical r2 blocker (Theorem C maximum coincides with
{[}15, Thm 1.1{]}) --- status in
v1.2}\label{g-historical-r2-blocker-theorem-c-maximum-coincides-with-15-thm-11--status-in-v12}

\begin{itemize}
\item
  \textbf{Now resolved (SV, manuscript text).} The citation is explicit
  in five places:

  \begin{itemize}
  \tightlist
  \item
    abstract l.25: "recovers the unrestricted extremal family of
    Okechukwu {[}15, Theorem 1.1{]}" and "No priority for the extremal
    value is claimed"
  \item
    §1 l.43, which lists all three clauses of {[}15, Thm 1.1{]} and Cor
    1.2 with "No priority is claimed"
  \item
    Theorem C l.84 (definition credited), l.95 (strengthening of {[}15,
    Thm 1.1{]}) and l.97 ("maximum and equality family are those of
    {[}15, Theorem 1.1{]}")
  \item
    §8.1 l.1245
  \item
    l.1014
  \end{itemize}

  \textbf{Blocker resolved: OK.}
\item
  \textbf{Logic of "Theorems C and C′ imply all three clauses" (l.1245)
  (AI).}

  \begin{itemize}
  \tightlist
  \item
    \emph{All-orders clause.} For n \textless{} N\_s, cp(G) ≤ c\_4(G) ≤
    e(G) ≤ C(N\_s,2), and Q\_s(n) ≥ −C(s+1,2) (Q\_s can be negative at
    tiny n for s ≥ 1). So K\_s = C(N\_s,2) + C(s+1,2) gives c\_4 ≤
    Q\_s(n) + K\_s at all orders. \textbf{Valid.}
  \item
    \emph{Eventual maximum clause.} Given by Theorem C and the witness.
    \textbf{Valid.}
  \item
    \emph{Equality clause.} Given by Cor 6.3 (6.18), which relies on
    Theorem C′ at δ = 0 plus the unrestricted lower witness, for n above
    max(N\_s\^{}stab, N\_s). \textbf{Valid as an eventual statement},
    which is also how {[}15{]} states it.
  \item
    \emph{Wording.} Saying "Theorem C′ recovers the equality
    classification" (l.43, l.109) is slightly compressed, since the
    proof needs Cor 6.3 and the Appendix E witness. \textbf{OK.}
  \end{itemize}
\item
  \textbf{Residual --- MINOR.} The abstract (l.25) and Theorem A (l.61)
  name only {[}5{]} as an alternative source for the all-orders Erdős-81
  bound. {[}15, Cor 1.2{]} also states it, in cp form, at every order,
  and {[}15{]}\textquotesingle s abstract says it answers the EOZ
  question. §1 l.43 does credit this. Adding "{[}15, Corollary 1.2{]}"
  next to "{[}5{]}" at l.25 and l.61 would remove any residual
  asymmetry.
\end{itemize}

\subsection{H. Overclaim scan}\label{h-overclaim-scan}

\begin{itemize}
\tightlist
\item
  \textbf{Universal b = 0.} Not claimed. It is explicitly disclaimed at
  l.27, l.795, l.1379 and l.2349. \textbf{OK}
\item
  \textbf{Universal linear mixed packing gap.} Not claimed. It is
  disclaimed at l.303 and l.2351. \textbf{OK}
\item
  \textbf{Small or effective threshold.} None claimed. The thresholds
  are called "explicit but impractical" (l.27, l.795, l.1126, l.2139).
  \textbf{OK}
\item
  \textbf{Unqualified superiority.} None found. The comparisons carry
  disclaimers (l.1215, l.1239, l.1245, l.1255, l.1267).
\item
  \textbf{Ambiguity at l.314 --- MINOR/UV.} The line says "The near
  threshold has been reduced from 10\^{}32 to 4·10\^{}12". It does not
  say whose 10\^{}32. {[}5{]} uses N = max\{10\^{}32, T(η/2)\} ({[}5{]}
  (1.3), p.2). The sentence may refer to an earlier version of the
  manuscript\textquotesingle s own threshold, but readers may read it as
  an improvement over {[}5{]}. Recommend naming the baseline.
\end{itemize}

\subsection{Additional antecedents noticed (informational, not cited; no
severity)}\label{additional-antecedents-noticed-informational-not-cited-no-severity}

\begin{itemize}
\tightlist
\item
  erdosproblems \#81 forum comments: Woett (19 Jun 2026), an explicit
  EOZ constant c ≥ 1/133; JamalAgbanwa (20 Jun 2026), a conditional n²/6
  + O(n) bound under a "weak Clique-Drop Hypothesis" (Zenodo 20778270,
  not opened).
\item
  Bo Ning, arXiv 2608.11536 / 2609.20305 (cp − cc difference): not
  related to chordal graphs.
\item
  No 2026 work was found, besides {[}5{]} and {[}15{]}, proving
  stability or c\_4 bounds for rooted simplicial defect. This is not a
  completeness claim.
\end{itemize}

\subsection{Summary of verdicts}\label{summary-of-verdicts}

{\def\LTcaptype{none} % do not increment counter
\begin{longtable}[]{@{}ll@{}}
\toprule\noalign{}
Item & Verdict \\
\midrule\noalign{}
\endhead
\bottomrule\noalign{}
\endlastfoot
A (Theorem C vs {[}15, Thm 1.1{]}; L-cutoff statement l.1255) & OK \\
A.3 (Table 4 / l.1257 description of {[}15{]}\textquotesingle s
mechanism) & MINOR (fix recommended) \\
B (cited items of {[}15{]}; order-3 scope) & OK \\
C (positioning vs Cor 5.5 / Thm 1.3) & OK \\
C.1 ("stronger hypothesis" wording, l.1259) & MINOR \\
C.2 ("order-four version of (1.5)", l.1261) & MINOR \\
D ({[}22{]} description and Lemma 11 case mapping) & OK \\
E ({[}5{]} claims, trust boundary, sections, m/9, A/2, credits;
unchanged HEAD) & OK \\
E (Anonymous attribution / Jacobian disclosure) & MINOR \\
F (Cipollini claim 201) & OK \\
G (r2 blocker) & Resolved (OK); residual abstract/Thm A asymmetry
MINOR \\
H (overclaims) & OK; l.314 baseline ambiguity MINOR/UV \\
\end{longtable}
}

No BLOCKER candidates and no MAJOR findings in E7.

\subsection{Coverage limits}\label{coverage-limits}

\begin{itemize}
\tightlist
\item
  \textbf{Databases not searched.} No MathSciNet, zbMATH Open or Google
  Scholar. dblp failed. There was no full-text search of 2026 journals,
  and "no later version / no journal version" rests on arXiv, Crossref
  and a general web search only.
\item
  \textbf{Linked sources not opened.} The Overleaf manuscript of
  {[}21{]} and Zenodo 20778270 were not opened.
\item
  \textbf{{[}5{]} scope.} Only README.md and
  lean/FORMALIZATION\_STATUS.md at cbde8a0a were read.
  docs/VERIFICATION\_STATUS.md, the Lean sources and the certificates
  were not.
\item
  \textbf{Depth of checks.} Proofs in {[}15{]}, {[}22{]} and {[}5{]}
  were not re-verified. The comparisons check statements, numbering and
  the logical role of lemmas. The analysis in A.2 is an inference about
  the literal chain, not an impossibility claim.
\item
  \textbf{Manuscript coverage.} Only the manuscript lines listed in the
  brief plus grep-located mentions of {[}5{]}, {[}15{]}, {[}21{]},
  {[}22{]} and {[}23{]} were read. Other appendices were not audited.
\end{itemize}

\begin{center}\rule{0.5\linewidth}{0.5pt}\end{center}

\subsection{E7 --- Lead-auditor verification of the delegated first
pass}\label{e7--lead-auditor-verification-of-the-delegated-first-pass}

The delegated first pass (\texttt{E7\_FIRST\_PASS.md}, same model
family) was spot-checked by the lead auditor against the saved literal
sources \texttt{sources/okechukwu\_v1.txt} (sha256 feff7419\ldots,
derived from PDF 9654af66\ldots) and \texttt{sources/verclos\_v1.txt}.

\subsection{Independently confirmed from the literal text of {[}15{]}
v1}\label{independently-confirmed-from-the-literal-text-of-15-v1}

\begin{itemize}
\tightlist
\item
  Definition (p.1--2): s-simplicial = d\_H(v) − ω(H{[}N\_H(v){]}) ≤ s;
  rsd(G) ≤ s: in every induced H and outside every proper clique (empty
  allowed) an s-simplicial vertex exists. Equivalent to
  \texttt{RootedDefectAt} (∃ clique C ⊆ N\_U(v) with
  \textbar N\_U(v)\textbar{} ≤ \textbar C\textbar{} + s ⇔ ω(N) ≥
  \textbar N\textbar{} − s). SOURCE\_VERIFIED.
\item
  Theorem 1.1 (p.2): constants K\_s, N\_s; cp(G) ≤ Q\_s(n) + K\_s at
  every order; for n ≥ N\_s the maximum is exactly Q\_s(n); equality
  exactly for (K\_\{k−s\} ⊔ K̄\_s) ∨ K̄\_\{n−k\}, k nearest to
  (2(n+s)+1)/6. Corollary 1.2: chordal case, cp ≤ F\_n + O(1) at every
  order, eventual max F\_n with classification. SOURCE\_VERIFIED.
\item
  Abstract of {[}15{]} states the chordal all-orders bound "answering a
  question of Erdős, Ordman and Zalcstein".
\item
  Theorem 1.3 (p.3): sequences, rsd = o(n), cp ≥ n²/6 − o(n²) ⇒ cliques
  A\_j of size n/3+o(n) with D\_A + e(G−A) = o(n²); and (1.5) cp(G) ≤
  (2n+1)²/24 + n·rsd(G) + ϵ(n)n² for every graph of order n ≥ 8.
\item
  Proof of Thm 1.3 (p.11): L fixed from Theorem 3.3 (L chosen from ε,
  p.9 "choose ξ \textgreater{} 0 and then L ≥ 4 \ldots"); "For (1.5),
  use L = 4 in Lemmas 3.1 and 3.4". Lemma 3.4 (p.10): q\_L ≤ q*\_L + ζn²
  (fixed L). SOURCE\_VERIFIED.
\item
  Proof of Thm 1.1 (p.19--20): minimal-counterexample argument via
  Proposition 5.4, whose final count (5.11) is produced by the
  edge-disjoint-triangle constructions of §4 (Lemma 4.1/4.2); Lemma 3.4
  enters only through Theorem 1.3\textquotesingle s localization.
  SOURCE\_VERIFIED (reading).
\item
  Corollary 5.5 (p.20): qualitative ε--δ edit stability for fixed s
  under cp ≥ Q\_s(n) − δn². §7 (p.23) asks for a quantitative version.
  SOURCE\_VERIFIED.
\item
  Only v1 listed on the arXiv abs page retrieved 2026-09-30 12:28 UTC.
\end{itemize}

Consequences for the manuscript: the v1.2 text now credits {[}15,
Theorem 1.1{]} for the class, Q\_s, the eventual unrestricted maximum
and the equality family (abstract l.25, §1 l.43, Thm C l.84--97, Cor 6.3
l.1014, §8.1 l.1245). \textbf{The historical r2 blocker F-01 is resolved
in this target} (it remains a historical finding for v1.0.1). The c4
upper bound is a genuine statement-level strengthening; neither the
manuscript nor this audit asserts that {[}15{]}\textquotesingle s method
cannot yield it.

\subsection{Confirmed MINOR items from the delegated pass (lead auditor
agrees)}\label{confirmed-minor-items-from-the-delegated-pass-lead-auditor-agrees}

\begin{itemize}
\tightlist
\item
  E7-M1 Table 4 (l.1210) and l.1257: describing
  {[}15{]}\textquotesingle s use of K\_3\ldots K\_L as a "far transfer"
  branch that "pays the budget \ldots{} in the branch with the
  corresponding margin" does not match {[}15{]}: Lemma 3.4 is used for
  localization (Thm 1.3) and for (1.5); Thm 1.1\textquotesingle s
  partitions come from the §4--5 triangle constructions.
\item
  E7-M2 l.1259 "stronger near-extremality hypothesis on c4": logically
  the weaker hypothesis (see E1-O3).
\item
  E7-M3 l.1261 "order-four version of its approximation (1.5)":
  {[}15{]}\textquotesingle s own proof of (1.5) already uses L = 4; only
  its stated conclusion is for cp. Wording could mislead
  (OBSERVATION/MINOR).
\item
  E7-M4 Abstract l.25 and Thm A l.61 name only {[}5{]} as an alternative
  source of the all-orders chordal bound; {[}15, Cor 1.2{]} (and
  {[}15{]}\textquotesingle s abstract) also claim it. §1 l.43 does
  credit {[}15, Cor 1.2{]}, so this is MINOR.
\item
  E7-M5 {[}5{]} authorship shown as "Anonymous"; the same group authored
  the Jacobian tool {[}14{]} used by the author. Transparent policy, but
  a one-line provenance note would help (MINOR; same as r2 E7-C6).
\item
  E7-M6 l.314 "reduced from 10\^{}32 to 4·10\^{}12": the prior
  value\textquotesingle s owner is unstated; {[}5{]} also uses 10\^{}32
  (per the delegated pass; not re-verified by the lead auditor) ---
  clarity.
\end{itemize}

\subsection{{[}22{]} spot checks}\label{22-spot-checks}

\begin{itemize}
\tightlist
\item
  Lemma 3 (p.3): n ≥ ε\^{}\{-1\}, G{[}Y{]} chordal, p\_G(v) ≤ εn² on X ⇒
  6ε\^{}\{1/2\}-close to chordal. Manuscript (F.10) states constant 8
  and no n ≥ 1/ε; comparison accurate. SOURCE\_VERIFIED.
\item
  Lemma 11 (p.10): tree T with at most k leaves, G0 of size ≤ k, subsets
  sampled uniformly of size m11 = 2\^{}40 ε\^{}\{-4\} k\^{}6 ln\^{}4 k.
  The manuscript\textquotesingle s adaptation (sequences with
  replacement, K nodes, enlarged pinned class, m11 =
  ⌈2\^{}58(K+1)\^{}15/ε\^{}12⌉) is described as an adaptation, not a
  literal statement. Accurate.
\item
  Only v1 on arXiv (abs page 2026-09-30).
\end{itemize}

\textbf{E7 first-pass status: PASS with MINOR findings (no attribution
blocker).} Coverage limits: arXiv API, Crossref, GitHub API,
erdosproblems.com, one general web search; no MathSciNet/zbMATH/Scholar;
Overleaf manuscript of {[}21{]} not opened; {[}5{]}\textquotesingle s
Lean sources not inspected. Not a complete novelty search.

\end{document}
